\RequirePackage[T1]{fontenc}
\documentclass[conference,10pt]{IEEEtran}
\usepackage{amsmath,amssymb,graphicx,url}
\title{Current-STA Transactional Repair for Timing--Resource Tradeoffs in 3D FPGA Routing}
\author{\IEEEauthorblockN{Anonymous development manuscript}}
\begin{document}
\maketitle
\begin{abstract}
Legal FPGA routes occupy distinct timing--resource points. In a fixed 3D graph, vertical sharing may save occupancy yet restrict critical connections or create capacity competition. Our current-STA repair layer searches complete connection groups on private route-tree copies and commits after physical, capacity and native timing checks. Retained-crossing domains encourage sharing; failed searches release only that network's domain while preserving other restorations. Three automatic repairs improve a strong native BNN route from 13.607436 to 13.228146 ns with forty additional planar wire units. Complete resource-first timing repair reaches 13.228146 ns with 43 SB and 5060 CB, versus 82 SB and 5062 CB for its native control, at a cost of 1762 more wire units. A fixed-terminal cut certifies minimal total vertical occupancy at this point. An independent five-connection feedback intervention improves ten setup components without regressions. On eltwise, a uniform K8 grouping cap reduces actual SB occupancy to 24 and CB occupancy to 965 at unchanged delay; an R1-based SPMV repair remains dominated by strong A0. A fixed-seed R1 resource trajectory reaches 33 SB at unchanged 13.421610 ns, spending 1721 more wire units. A local budgeted resource/timing composition reaches 13.228146 ns with 47 SB, trading four extra SB for 313 fewer wire units than the 43-SB point. Matched controls distinguish selection, grouping and domain contributions. Fresh-route transfer also reduces softmax SB from six to three under both controls, while SPMV has no eligible cut. On eltwise, acceptance-ordered execution preserves committed states while reducing observed complete cost from 289.21 to 203.48 s; complete pooled-strategy timing and broader fixed-policy confirmation remain open.
\end{abstract}
\begin{IEEEkeywords}
3D FPGA, detailed routing, static timing analysis, transactional repair, multiobjective optimization
\end{IEEEkeywords}

\section{Problem and Motivation}
Three-dimensional FPGA fabrics couple directed reachability, scarce vertical channels and detailed timing. A different route can improve timing, save resources or exchange one benefit for another even when packing, placement and architecture remain fixed. A repair that does not reduce the final critical-path delay (CPD) may still be useful under an explicit resource or timing requirement. Conversely, a locally shorter connection is not evidence of a better complete route: another endpoint can become critical, or the replacement may prevent a different network from using a scarce channel.

VTR provides the detailed routing and timing infrastructure used here~\cite{vtr9paper}. The placement and architectural context follows the separation of design tradeoffs emphasized by Escaping Flatland~\cite{flatland}; our comparisons hold each physical input fixed. This work concerns post-legal routing decisions. It neither changes the physical timing model nor infers silicon area or power from reduced occupancy on a fixed device.

Shared-crossing proxy scores can diverge from final timing through candidate coverage, retained branches and native convergence. We therefore select groups using current STA, search native route-tree replacements and validate actual outcomes. The intervention comprises group selection, dependency scope, retained domains, acceptance and scheduling.

We distinguish legal quality tradeoffs, matched mechanism controls and complete configuration cost. Serial full-flow measurements complement development evidence; broader frozen-policy and pooled measurements remain required.

\section{Objective and State Identity}
An input identity fixes the directed resource graph, packed netlist, placement, endpoint mapping, SDC and clock domains. A legal route $R$ uses actual graph edges and respects node capacities and native routing constraints. We record
\begin{equation}
 q(R)=(D(R),W(R),S(R),C(R),T(R)),
\end{equation}
where $D$ is native final CPD, $W$ is planar wire in the auditor's native resource units, $S$ is occupied SB pairs, $C$ is occupied CB output sources, and $T$ is comparable complete algorithm time. Paired SB half-nodes are one vertical resource; multiple CB edges from one output are not independent outputs. Guide recommendations are not occupancy measurements.

A final legal candidate is retained when it is not dominated on the measured, comparable coordinates. Quality projections use the first four coordinates and preserve uncertainty from printed CPD. They do not establish five-dimensional dominance when complete time is unknown. Different convergence rules, repair opportunities and implementation versions remain explicit, even for a descriptive comparison on the same physical inputs.

Final retention is separate from a branch's transition rule. We retain a strict CPD mode, $D(R')<D(R)$, and a fixed-resource mode that reduces CPD within registered absolute resource limits. A resource-under-timing mode permits wire tradeoffs within a fixed envelope while requiring nonincreasing selected resource objectives and at least one strict reduction. The current resource probes select $(S,C)$ as those objectives, cap CPD at the root value, cap wire at the root value plus 0.1\%, and prohibit increases beyond root vertical counts. These are design parameters, not general termination conditions.

An endpoint-vector control fixes all endpoint/launch/capture identities and compares complete native setup slack values. It accepts only nonworse CPD, no setup-component regression and at least one strict setup improvement, with no increased hold violation. Native numeric values are compared without an epsilon. Resource policies impose their stated CPD and hold conditions; they do not silently inherit the stricter endpoint-vector requirement. The new bounded vector variant additionally enforces the strict timing control's immutable wire/SB/CB limits. Anchors do not roll forward, and no tolerance can accumulate timing degradation.

\section{Mechanism: Current-State Whole Transactions}
\subsection{Select targets and dependency scopes}
Timing selection follows current native winning paths, including tied fan-ins. Capacity queries identify blocking owners; grouped targets and owner scopes release dependencies frozen by single-connection repair.

Resource selection scans the actual current route trees. A candidate removes all sinks below a selected crossing subtree, preserves paths to the other sinks, and requires both structural release potential and a retained crossing alternative. The screen uses current slack margin; it is not a proof of final resource savings. Ranking is deterministic: released crossings per removed sink, fewer sinks, larger current margin, then network and pin IDs. A K8 group contains at most eight distinct networks; fewer are used when fewer qualify. Removed connections are restored in descending current raw criticality with deterministic ties. Accepted states receive a fresh current selector.

\subsection{Native searches and retained-crossing domains}
For each group we copy the selected trees and prune exactly the registered sink sets. Occupancy is updated as every local replacement is installed, so later searches see earlier resource use. Native weighted search uses the full device search box and the existing timing/congestion objective. Capacity-infeasible results are not committed. Real switch delays and final STA remain unchanged.

The domain control initially allows cross-layer edges only through crossing sources retained after pruning. It preserves useful sharing but can become infeasible under current capacities. The unrestricted native control receives the same current source, sink set, tree scope, objective and search range. Paired development probes evaluate both variants and use the registered domain-first acceptance order; native-only trajectories instead reselect from their own accepted unrestricted states. Those trajectories must not borrow a guided successor after divergence. Resource-policy execution can instead try the preferred domain candidate first and omit the optional native diagnostic only after complete resource acceptance. A rejected or infeasible domain candidate retains the same-source native fallback; omitted candidates have no inferred metrics.

Automatic failure release makes the domain a revisable decision. If guided native search fails to return a capacity-feasible branch, search costs are reset and only that network's domain is removed. Previously successful restorations and their occupancy stay in the private transaction. Other networks remain guided. Search failure is an operational trigger, not itself a mathematical infeasibility certificate. A separate capacity diagnosis can justify a stronger statement for a particular fixed scope.

\subsection{Complete checks, acceptance and restoration}
All selected connections must be restored before acceptance. Independent checks validate real edges, unchanged retained branches and terminals, actual global capacities and resource counts. Native FULL and fresh STA evaluate the actual replacement. Acceptance commits the entire group; failure or rejection retains the source state. Intermediate private states may spend resources or slack when a complete mechanism-backed combination is subsequently checked. Requiring every component to improve CPD would exclude some useful groups.

Candidates receive native LOAD/FULL/STA, or are installed in a private worker for FULL/fresh STA and dedicated export. The latter restores trees, sink lookups, occupancy, digest and hooks before source reobservation. Tiny, BNN-resource and free-path timing tests match separate same-binary LOAD/STA plus physical/scope checks; large forced-path timing transactions remain unvalidated.

Loaded post-legal contexts reconstruct occupancy and native STA but reset the original route run's accumulated PathFinder price history. Matched native/domain controls share this context. An eventual in-router integration must explicitly preserve or reset that history and evaluate the choice.

\section{Overall Development Results}
The implementation derives from VTR revision \texttt{0b912334fdd83bb754b302146ac0d1ec9edc9a0e}. Fixed-input serial clustered routing uses two-layer edge-encoded Hybrid-O/SB-only/CB-only graphs, ideal clocks and native FULL/STA. Each record binds physical inputs, constraints, binary, request and acceptance. Flat routing, CHANZ and general multilayer remain unvalidated. A0 is native; R1 is the older guided method. Strong BNN roots permit eight native convergence attempts with a 1.01 threshold. BNN has no supplied SDC: defaults use one ideal ap\_clk domain and zero input/output/setup/hold constraints; bound echoes fix these. U/Tf remain separate versions/opportunities.

\begin{table}[t]
\centering\caption{BNN legal quality points. Fast A0 follows three current-STA commits from pristine strong A0. Resource inputs/opportunities remain separate; both 47-SB local timing menus are terminal. Quality comparison, not matched-compute ranking.}
\label{tab:bnn}
\small
\begin{tabular}{lrrrr}\hline
State & CPD(ns) & Wire & SB & CB\\\hline
Strong A0 &13.607436&2059567&94&5062\\
Strong R1 &13.421610&2068908&86&5072\\
A0 native K8 (end) &13.607436&2059591&82&5062\\
A0 guided K8 (end) &13.607436&2061353&43&5060\\
Fast A0/native K8 &13.228146&2059603&86&5062\\
Fast A0/guided K8 &13.228146&2059882&79&5061\\
Fast R1/guided K8 &13.228146&2069040&74&5071\\
Full native + timing &13.228146&2059627&82&5062\\
Full guided + timing &13.228146&2061389&43&5060\\
45 SB vector feedback &13.599798&2061253&45&5060\\
45 SB vector-rejected &13.413972&2061257&45&5060\\
R1 selective release &13.228146&2069104&72&5071\\
R1 fixed-seed K8 &13.421610&2070629&33&5071\\
Budget64 (local) &13.607436&2061036&47&5060\\
47 SB local native end &13.228146&2061249&47&5060\\
Budget64 local end &13.228146&2061076&47&5060\\\hline
\end{tabular}
\end{table}
The strong-A0 timing trajectory accepts three current-state transactions and reduces CPD by 2.787376\%, adding forty wire units with SB/CB unchanged. Its third step improves CPD but regresses a noncritical setup component, making the distinction between strict CPD and endpoint-vector policies observable. The same-source fixed-vertical policy independently selects the identical three states using exact completed-query/transaction reuse; its replay time is not evidence of speedup.

From fast A0, native K8 ends at 86 SB/2059603 wire; guided K8 spends 279 more wire for 79 SB/5061 CB. Conversely, timing after a historical two-accept resource-first 79-SB prefix reaches 13.228146 ns under those vertical caps and the original wire envelope. Its final nine-network/twelve-connection exchange spends 1.114955 ns of noncritical margin for a 0.185826 ns critical gain. Strict CPD accepts while the vector control rejects; the phase has eleven setup improvements, seventeen regressions and no hold-violation increase. Final quality/vectors match timing-first/guided-K8 despite different ordered bytes: convergence to a known point, not an order-commutation theorem. Complete chains below supersede these historical prefixes as native bounds. All ten winning switch distances equal optimistic source-only references; occupancy/shared-tree/native-arithmetic relaxations prevent a global native-STA bound.

\begin{table}[t]
\centering\caption{Eltwise fixed-seed K8 controls from the same pristine A0 root, using uniform grouping and acceptance. Both branches use their own successors.}
\label{tab:eltwise}
\small
\begin{tabular}{lrrrr}\hline
State & CPD(ns) & Wire & SB & CB\\\hline
Strong A0 root &6.018101&217065&30&968\\
Native K8 &6.018101&217050&26&968\\
Paired fixed-seed K8 &6.018101&217183&24&965\\\hline
\end{tabular}
\end{table}
Eltwise guided K8 uses actual group sizes six then three. It saves one additional SB relative to the prior guided K4 endpoint, with identical CPD, wire and CB. No failure release triggers in this circuit, so the gain is attributed to grouping/domain decisions rather than fallback. The native-only branch accepts one move, then rejects three legal but resource-worse alternatives; these rejected candidates remain recorded. On SPMV, paired and native-first R1 repair produce the same legal route at 7.045986461 ns,148605 wire,zero SB and 967 CB, but remain quality-dominated by strong A0. Initial routability and a circuit-balanced timing improvement across the suite have not been established by these repair probes.

\section{Pareto Tradeoffs and Attribution}
\begin{figure*}[t]
\centering\includegraphics[width=.98\textwidth]{data/joint_quality_projection_24.pdf}
\caption{Archive 499 contains 305 BNN observations on one physical-input identity, including strong A0/R1/U/Tf, complete-resource timing, independent 45 SB feedback and all legal fixed-seed candidates. Duplicates collapse to 189 plotted coordinates. Both 47-SB timing menus are local terminals. Roots/convergence/opportunities remain explicit in data; projections omit other resource coordinates and complete time and are not two-dimensional Pareto boundaries.}
\label{fig:joint}
\end{figure*}
Table~\ref{tab:bnn} and Fig.~\ref{fig:joint} retain external strong baselines. The fast A0 native-K8 point quality-dominates the earlier slower A0 guided-K4 point, which had the same 86 SB and 5062 CB but wire 2059667. Slower native-first low-wire states and the CB 5059 reassignment state retain distinct resource tradeoffs; the new timing prefix quality-dominates several CB 5060 predecessors. Quality domination is not sufficient for a complete-runtime claim when comparable time is missing.

\subsection{Extending the native control changes the comparison}
The completed native-order resource cohort ends at 13.607436067 ns for both own-current branches. The unrestricted control makes five accepts, using 2059591 wire,82 SB and 5062 CB. The paired domain/failure-release branch makes nine accepts, using 2061353 wire,43 SB and 5060 CB. Its additional 1762 wire units buy 39 fewer SB and two fewer CB. Both exhaust their registered current menus; historical two-accept points are prefixes, not native resource bounds. Native and paired workflows cost 2061.28 and 7293.64 s and create nine and nineteen native children respectively. The paired branch pays its additional searches; these imported-state measurements are not equal-compute or cold whole-configuration timing.

Continuing from each completed resource endpoint, three strict current-STA commits reach the identical 13.228145690 ns CPD. The native branch uses 2059627 wire/82 SB/5062 CB; the paired branch uses 2061389/43/5060. The extra 1762 wire units still buy 39 SB and two CB, now at the faster timing point. Both are nondominated in the collected four-axis population and both timing menus terminate. Full-vector comparison has 550 setup improvements and 799 regressions, so equal CPD does not imply endpointwise dominance. The native imported timing phase costs 3093.90 s; the paired recovery segment costs 450.11 s, with original interrupted elapsed time unknown. Those durations are not comparable complete costs.

Fixed terminal layers require 5103 counted nonglobal networks to cross. Scanning every graph crossing verifies unit-capacity CB sources and uniquely paired unit-capacity SB halves, so distinct networks cannot share a counted resource. Thus $S+C\ge5103$, attained by 43/5060. Under $C\le5060$, $S\ge43$ is also attained. This proves fixed-graph vertical occupancy optimality, not wire, CPD or placement optimality.

A separate historical composition starts from the earlier seven-accept 45-SB/5060-CB prefix. Three strict repairs reach 13.228145690 ns at 2061301 wire/45 SB/5060 CB in 3028.24 s of imported repair. This prefix is not the final 43-SB resource input. Its full-vector branch accepts none of six proposals in 2747.93 s: one target improves physically but leaves the endpoint vector unchanged because another tied path remains; larger combinations worsen donors.

Same-source rejection feedback drops delay-increasing donors and jointly covers two tied-path targets with three donors. The five-connection transaction reaches 13.599797732 ns (7.638334 ps gain) at 2061253 wire/45 SB/5060 CB. All ten changed setup components improve, none regress, and no hold violation increases; strict/vector policies accept. Its 434.52 s intervention includes preparation, coordination and evaluation, retaining a distinct low-wire point. From its actual successor, a twelve-connection current-owner exchange reaches 13.413972155 ns at 2061257 wire/45 SB/5060 CB: 40 fewer wire than the earlier equal-timing/vertical point. One setup component improves and one regresses, remaining 2.487 ns above the new worst, without increased hold violation. Strict accepts; the registered vector branch rejects. Both exchange orders yield identical ordered routes/full vectors; 409.82 s excludes preparation. The controller ends thirteen proposals with only the known feedback commit and six rejected successor candidates. Its 3755.45 s imported execution counts seven native coordinators, the repeated primary and six historical reuses; initial routing/historical creation are excluded.

\subsection{Automatic transfer and deterministic repair seeds}
With delay/RR-ID seeds fixed globally, eltwise native/domain roles end at 217050/217183 wire,26/24 SB,968/965 CB and 6.018101129 ns in 123.75/126.10 s of imported work. BNN A0 ends at 2059591/2061377 wire,82/43 SB,5062/5060 CB and 13.607436067 ns in 2076.31/7099.32 s. Paired wire decreases by 8 on eltwise but increases by 24 on BNN versus native-order predecessors. Six of eight root/control roles are sealed on two circuits; roots are not independent circuit samples.

Serialization changes eltwise from 217214 to 217191 wire under identical requests/domains; reordering only net 6719 restores the old FULL/STA candidate. Six transactions verify explicit RR-ID/delay-RR-ID seeds yielding 217214/217183 wire, both 24 SB/965 CB at unchanged CPD. Changed later costs rule out heap ties alone. Eighteen regressions check defaults/restoration; transfer fixes delay/RR-ID.

The fixed policy's eight R1 commits reduce 86 SB/5072 CB to 33 SB/5071 CB at unchanged 13.421610490 ns, adding 1721 wire to 2070629. Its 37-SB intermediate and 33-SB final points are nondominated. At rounds 4--7, same-source unrestricted candidates fail the resource rule while guided candidates pass. The own-current native trajectory remains checkpointed, so whole-trajectory superiority is unproven. All eighteen legal candidates remain; 5046.61 s and nineteen native children describe imported work, not cold cost. Tf makes twelve commits from 105 SB/5062 CB to 43 SB/5060 CB at unchanged 13.421610490 ns, adding 1721 wire to 2076038. Its twenty-four legal outcomes and 6752.95 s imported cost are sealed, but the endpoint is dominated by strong-A0's 43-SB timing point; Tf-native closure remains pending.

\subsection{A connection budget exposes another resource tradeoff}
A variant caps current-ranked connections at 64 without splitting cuts; if none fits, the first indivisible cut runs alone. At one BNN 48-SB state it selects network 4111's 54 connections versus the original three-network/462-connection group. OFF adds one SB and improves 85 setup components without regressions. ON spends 40 wire, reaching 2061036 wire/47 SB/5060 CB at unchanged 13.607436067 ns; setup 57 improve/24 regress, with no hold violation increase. Resource acceptance passes, vector acceptance rejects. This nondominated point saves 177 wire versus the larger native 47-SB result at 2061213 wire; larger guided 45 SB/2061285 wire remains a tradeoff. Query cores 51.31/60.11 s versus historical 509.71/476.90 s are component observations. Shared imported-scope execution 948.00 s includes actual loads/checks but excludes preparation/initial routing; it is not a complete budgeted trajectory or full-flow speedup.

From the same guided 48-SB source, unrestricted 462-connection and budget64 guided 54-connection alternatives both complete three current-STA commits under the original 2061627-wire/47-SB/5060-CB envelope. Both end at 13.228145690 ns, but guided uses 2061076 wire versus native 2061249: 173 fewer at identical vertical counts, dominating this local control. Against the complete 43-SB point it saves 313 wire for four more SB. Guided's 13.413972155-ns/2061072-wire intermediate remains a tradeoff.

All three timing requests, connection-delay changes and endpoint changes match. The first replacement's complete physical trees agree, preserving 177 wire; later trees differ, shrinking that advantage to 173. These are different full source states, not byte-equivalent trajectories or one-factor domain/grouping isolation. Actual complete timing segments cost native/guided 3056.67/3008.02 s, with source reobservation 210.01/210.48 s separately. Resource-prefix creation/initial routing are excluded; neither is a pristine whole-native trajectory or cold complete algorithm result.

\subsection{Capacity feedback changes the domain decision}
An all-eight-domain R1 transaction fails at network 31767/pin 1 after five successful restorations. A directed path exists when occupancy is ignored. Under the recorded transient occupancy, the retained domain has no capacity path while an unrestricted witness exists. Reconstructing the minimum occupancy immediately after pruning also leaves no domain-capacity path. Because subsequent restorations only add occupancy, changing order or search pressure cannot fix this particular frozen scope/domain. The diagnostic searches a relaxed superset of native re-expansion/closure choices; its negative result is a necessary infeasibility certificate for that scope, not for the entire FPGA.

Releasing only that network's domain yields a fully legal state at 72 SB, compared with 74 SB for the same-source unrestricted control, at the cost of 64 wire units and unchanged CPD/CB. The automatic native release implementation obtains exactly the prior selective-domain route and full endpoint vectors without a manually supplied release network or restarting earlier restorations. Nine setup components improve and none regress, so endpoint-vector acceptance succeeds while strict CPD rejects. All 23 components in a fixed source near-critical band stay unchanged. This step demonstrates resource and secondary-timing benefit, not later CPD improvement.

\subsection{Grouping contribution}
Fast-A0 native K8 matches two K4 steps in routes/vectors. Its first K4 component is endpoint-flat and rejected, whereas K8 improves one setup component without regressions: evidence for whole combinations, not a pristine-root endpoint trajectory.

\section{Overhead, Limitations and Next Confirmation}
Table~\ref{tab:completecost} reports full current-binary configurations from supplied RR graph, architecture, BLIF, packed netlist and placement through graph-index construction, fresh routing and complete K8/rank 3 repair. Each circuit's initial route bodies and full vectors match across controls; one Census binary performs routing, repair and candidate STA. On eltwise, native accepts once in four scopes and domain-first twice in two. Pairing spends 133 wire for two fewer SB and three fewer CB; fixed seeds save 31 wire relative to native traversal.
\begin{table}[t]
\centering\footnotesize\setlength{\tabcolsep}{3pt}
\caption{Complete serial configurations. N/D: native-only/domain-first. Seeds are delay/RR-ID except $^\dagger$native traversal. Fused/lazy denote domain-first execution variants. Alternate-topology rows use eltwise; other rows use Hybrid-O. PSS is sampled process-tree sum in GiB.}
\label{tab:completecost}
\begin{tabular}{lrrrr}\hline
Circuit/role & CPD(ns) & Wire/SB/CB & Time(s) & PSS\\\hline
Eltwise N&6.018101&217050/26/968&290.19&1.274\\
Eltwise D&6.018101&217183/24/965&289.21&1.302\\
Eltwise D$^\dagger$&6.018101&217214/24/965&289.66&1.297\\
SPMV N&7.034698&148463/0/967&91.69&1.010\\
SPMV D&7.034698&148463/0/967&91.55&1.034\\
Softmax N&8.975747&169486/3/713&274.22&2.502\\
Softmax D&8.975747&169486/3/713&336.68&2.534\\
Eltwise fused&6.018101&217183/24/965&313.74&1.296\\
Eltwise lazy&6.018101&217183/24/965&203.48&1.310\\
SB-only N&6.040485&227439/999/0&424.59&1.271\\
SB-only D&6.040485&227459/999/0&533.57&1.282\\
CB-only N&5.862299&212511/0/982&74.15&0.740\\
CB-only D&5.862299&212511/0/982&72.46&0.740\\\hline
\end{tabular}
\end{table}

The fixed policy also completes both roles on SPMV and softmax. SPMV has no eligible reusable-crossing cut: both flows retain 148463 wire/0 SB/967 CB with zero repair, so there is no domain intervention or inferred global optimality. On softmax, both controls commit once, reducing 6 SB to 3 and wire 169490 to 169486 at unchanged 8.975747257 ns and 713 CB. Domain-first pays two native transactions versus one, adding 62.46 s in this observation without additional quality benefit. No failure release triggers in these three circuits. These known-circuit transfer checks support resource-refinement scope, not universal domain gains or unseen-test confirmation.

Fused same-process candidate FULL/STA preserves all four candidate vectors, ordered routes, decisions and final quality, but costs 313.74 s versus 289.21 s for separate LOAD/STA. Its two source-serialization controls cost 41.90 s inside the full command; fewer native processes do not establish speedup. Acceptance-ordered execution instead computes ON first and skips two OFF diagnostics only after legal resource acceptance. Every executed candidate and selected state matches exactly; queries fall from 64 to 32 and attempts from 101 to 51. Complete elapsed time is 203.48 s, 85.74 s (29.65\%) less in this observation. Historical matched controls remain preserved; skipped outputs are not imputed. The fixed schedule transfers to softmax with exact executed states: 271.70 versus 336.68 s, skipping one OFF. SPMV has no eligible cut or skipped work and takes 93.15 versus 91.55 s. These are single development observations, not statistical speedups.

A fixed-policy architecture transfer holds eltwise packing/placement and the audited n40/w100 plane constant, changing vertical connectivity. SB-only native/domain both reduce 1012 to 999 SB at unchanged 6.040485445 ns; domain ends 20 wire higher and costs 108.98 s more in this observation. One actual ON candidate reaches 996 SB but 6.114032391 ns, exceeding the fixed CPD cap; actual OFF fallback preserves 6.040485445 ns at 999 SB. This verifies policy fallback after full timing rejection, not a per-net search-failure release. CB-only has no eligible cut: its 982 counted fixed-terminal cross-layer nets require 982 distinct unit-capacity OPIN crossing sources, matching observed occupancy. This vertical lower bound does not bound wire or CPD. All four roles complete legally; domain benefit is architecture-dependent.

All commands count accepted/rejected candidates, FULL/STA, physical checks, restoration and final gates, without external route/query/candidate reuse. GNU RSS high-water spans 0.701--2.241 GiB. Half-second tree sampling may miss peaks; missing reads total sixteen across fifteen configurations and remain explicit per run. Packing/placement generation is upstream. Fifteen configurations cover three circuits and three vertical-connectivity families, with one shared-host observation each and uncontrolled OS cache; timings are descriptive, not statistical speedups.

Legacy serial native-order configurations use two bound binaries without failure release. Native/paired build costs 288.49/292.47 s belong to 217050/217214 wire and corresponding vertical counts. Graph-index reuse takes 221.11/226.35 s with exact initial/intermediate/final vectors, ordered routes, requests and decisions, including rejects. Observed decreases of 67.38/66.12 s exclude supplied index construction, separately 62.09/61.48 s. These single observations do not belong to the new 217183-wire output.

Timing gains concentrate in BNN; resource transfer does not establish initial routability. K1/K4/K8 and scheduling pilots remain distinct. Confirmation must freeze rules, extend circuits/seeds/architectures, and match A0/R1/native scope, targets, domains and acceptance opportunities; pooled complete time remains unknown.

\subsection{Validation overhead}
Streaming pin parsing preserves full vectors while reducing component time/memory; detailed measurements remain supplemental. Three completed native-count roles on two circuits preserve all sixteen legal candidates, including rejects, in full vectors, ordered routes, requests/domains and decisions. All counter fields match independently audited references; assurance remains native, and each selected final passes its own independent audit/context seal. Direct/native-count imported seconds are 101.02/88.87, 100.33/88.48 and 1804.56/1418.13. BNN-paired and cold pooled costs remain pending; these are not statistical speedups.

\subsection{Direct audit avoids another candidate serialization}
Direct auditing checks candidate identity, restoration, terminals, retained edges, untouched networks and capacity without another serialization. Six transactions and six corruption guards pass; domain/FULLSTA gates remain. Four roles on two circuits preserve all thirty-four legal candidates, including rejects: complete vectors, ordered route bodies, requests, domains and decisions match. Old/direct imported seconds are 123.75/101.02 and 126.10/100.33 for eltwise native/paired, and 2076.31/1804.56 and 7099.32/6307.79 for BNN native/paired. These exclude fresh routing; single observations are not statistical speedups. Forced reservations remain untested.

Reproduction starts from the source and input bindings in \path{multiobjective_budget47_complete_pair_01/completion.json}; its archive preserves all 499 observations, including dominated and failed-policy legal candidates. Projection/feedback receipts distinguish final states, interrupted costs and historical prefixes.

\section{Conclusion}
Current-STA whole transactions expose timing--resource alternatives beyond strict CPD retention. The complete 43-SB/5060-CB point attains the fixed-terminal vertical cut bound, trading 1762 extra wire for 39 fewer SB and two fewer CB at equal CPD versus native. Tied-path feedback also improves a full endpoint vector on a distinct input. R1-rooted refinement extends the observed low-vertical frontier to 33 SB. Matched native controls separate selection and grouping from retained-domain contributions. Acceptance-ordered execution removes unused control work while retaining exact selected states; broader frozen-policy confirmation and complete pooled-runtime measurements remain required.

\begin{thebibliography}{2}
\bibitem{vtr9paper} M. A. Elgammal et al., ``VTR 9: Open-source CAD for fabric and beyond FPGA architecture exploration,'' \emph{ACM Trans. Reconfigurable Technol. Syst.}, 2025, doi:10.1145/3734798.
\bibitem{flatland} C. Hao, A. B. Kahng, B. Pramanik, and I. Youssef, ``Escaping Flatland: A placement flow for enabling 3D FPGAs,'' in \emph{Proc. ACM/IEEE DAC}, 2026, doi:10.1145/3770743.3804424.
\end{thebibliography}
\end{document}
